\documentclass[12pt]{book}
\usepackage{amsmath}
\usepackage{graphicx}
\begin{document}
\chapter{Energy in Our Community}
Energy makes change possible. A moving bicycle has \emph{kinetic energy}. A book resting on a high shelf has gravitational potential energy because of its position.

\section{Where does electricity come from?}
Electricity can be generated in several ways. Solar panels convert light into electrical energy. Wind turbines transfer energy from moving air to a generator. Hydroelectric stations use moving water. Each method has benefits and limits.

\section{Investigate a school building}
Choose three devices in a classroom. Record the device, its power rating in watts, and the estimated hours of use each day. Energy use can be estimated with the equation:
\[
E = P \times t
\]
Here, $E$ is energy in watt-hours, $P$ is power in watts, and $t$ is time in hours.

\section{Compare the data}
\begin{table}[h]
\caption{Example classroom devices}
\begin{tabular}{lrr}
Device & Power (W) & Hours per day \\
LED lamp & 8 & 5 \\
Laptop & 45 & 4 \\
Projector & 220 & 2
\end{tabular}
\end{table}
Calculate the daily energy estimate for each device. Explain why the highest-power device does not always use the most energy.

\begin{figure}[h]
\includegraphics[width=0.75\textwidth]{energy-flow.pdf}
\caption{Figure 3.1. Sunlight reaches a solar panel; electrical energy flows from the panel to a classroom lamp.}
\end{figure}

\begin{enumerate}
\item Identify the energy source.
\item Describe one transfer or conversion.
\item State one limitation of the system.
\end{enumerate}
\end{document}
